% =====================================================================
% Informe de la semana --- nodo de adquisición de geófono
% Se compila con:  lualatex informe.tex   (dos veces, por las referencias)
% NO con pdflatex: usa fontspec.
% o con el script:  compilar.ps1
%
% OJO con los guiones: la raya (U+2014) y el guion medio (U+2013) NO existen
% en ec-lmr10, que es lo que toma lualatex con fontenc T1, y desaparecen del PDF
% sin mas aviso que una linea "Missing character" en el log. Se escriben --- y
% -- , que es lo que LaTeX espera.
%
% Las figuras las genera src/interfaces/python/figuras_semana.py a partir
% de los JSON crudos. Ningún número de este documento está escrito a mano
% sin que exista el archivo de datos que lo respalda; la sección final
% dice dónde está cada uno.
% =====================================================================
\documentclass[11pt,a4paper]{article}

% LUALATEX YA ES UTF-8 NATIVO, asi que inputenc/fontenc no van: con ellos
% los caracteres de un byte alto se remapean mal y el PDF sale con '¿'
% convertido en '£' y '×' en 'Œ'. Se ve en el PDF y NO en el log, que es
% lo peor: compila limpio y sale mal.
\usepackage{fontspec}
\setmainfont{Latin Modern Roman}
\setsansfont{Latin Modern Sans}
\setmonofont{Latin Modern Mono}
\usepackage[spanish,es-noquoting,es-nolists,es-noshorthands]{babel}
\usepackage[margin=2.4cm]{geometry}
\usepackage{graphicx}
\usepackage{booktabs}
\usepackage{longtable}
\usepackage{array}
\usepackage{seqsplit}
\usepackage{xcolor}
\usepackage{amsmath}
\usepackage[colorlinks=true,linkcolor=azuloscuro,urlcolor=azuloscuro]{hyperref}

\definecolor{azuloscuro}{HTML}{1F4E79}
\definecolor{rojo}{HTML}{C0392B}
\definecolor{verde}{HTML}{1E8449}
\definecolor{grisclaro}{HTML}{F2F2F2}

\newcommand{\si}{\textcolor{verde}{\textbf{sí}}}
\newcommand{\no}{\textcolor{rojo}{\textbf{no}}}
% Las rutas son una sola palabra y no rompen: se salen de la columna.
% Este \ruta deja cortar despues de cada barra, que es donde uno lo
% esperaria, y no en cualquier lado.
\newcommand{\ruta}[1]{{\ttfamily\small\seqsplit{#1}}}

\title{\vspace{-1.5cm}\textbf{Nodo de adquisición de geófono}\\[2mm]
       \large Resultados de la semana del 31 de agosto al 6 de septiembre de 2026; cierre de EXP4c el 7}
\author{}
\date{}

% Los numeros de la deriva y el ruido los escribe figuras_semana.py desde
% el mismo JSON que las figuras. EXP4c se cerro el 2026-09-07 a las 08:01.
% GENERADO POR figuras_semana.py. No editar a mano.
% Sale del mismo JSON que las figuras; se actualiza al recompilar.
\newcommand{\derivaArranco}{-14}
\newcommand{\derivaExtremo}{-144}
\newcommand{\derivaExtremoH}{3,7}
\newcommand{\derivaAhora}{-37}
\newcommand{\derivaNeto}{-22}
\newcommand{\derivaBandas}{1,3 h: 53 mV. 2,6 h: 123 mV. 3,9 h: 130 mV. 5,2 h: 130 mV. 6,5 h: 150 mV.}
\newcommand{\derivaCrecio}{21}
\newcommand{\derivaVeces}{4,4}
\newcommand{\derivaTramos}{-5,3, -20,5, -25,1, +9,0, +35,8}
\newcommand{\derivaRecta}{8}
\newcommand{\derivaHoras}{6,5}
\newcommand{\derivaMuestras}{367}
\newcommand{\derivaBanda}{150}
\newcommand{\derivaTasa}{1}
\newcommand{\derivaResiduo}{29}
\newcommand{\derivaCome}{26,2}
\newcommand{\ruidoN}{74}
\newcommand{\ruidoRms}{501}
\newcommand{\ruidoHz}{31}


\begin{document}
\maketitle
\vspace{-1.2cm}

% =====================================================================
\section*{Lo que hay que saber si se lee una sola página}

\textbf{El nodo se calibra solo.} Arrancando con las cuatro referencias en cero
---el estado de un nodo recién encendido en el campo--- y sin computadora conectada,
lleva el tap del filtro pasabajos desde fuera del rango medible hasta
\textbf{34 mV de $V_{ref}$}, en 579~s. Es la primera vez que lo hace.

El apartado de la tesis sobre autocalibración queda \textbf{completo}: las cuatro
preguntas que lo sostienen tienen ahora respuesta medida.

\begin{center}
\begin{tabular}{p{9.2cm} c p{4.4cm}}
\toprule
\textbf{pregunta} & & \textbf{evidencia} \\
\midrule
¿Sin calibrar el nodo funciona? & \no &
  12 de 14 combinaciones de ganancia arrancan con una etapa que no transmite \\[1mm]
¿Alcanza con calibrar una vez en fábrica? & \no &
  el punto recorre \derivaVeces{} veces la corrección y presenta tendencia lenta \\[1mm]
¿El nodo puede calibrarse solo? & \si &
  de fuera de rango a 34 mV de $V_{ref}$, en 579 s, sin PC \\[1mm]
¿Calibrar empeora el ruido? & \no &
  el RMS no crece a lo largo de \derivaHoras{} h \\
\bottomrule
\end{tabular}
\end{center}

\vspace{3mm}
\noindent\textbf{El hallazgo que lo destrabó todo} fue que el firmware se
detenía \emph{63 códigos antes} del punto donde la cadena empieza a responder,
por un límite simétrico que nadie había pedido. La figura~\ref{fig:adder} lo
muestra de una mirada.

\vspace{2mm}
\noindent\textbf{Y el error de método que más costó} fue mío: barrí un actuador
hasta el límite que le ponía el propio firmware y leí el resultado como una
propiedad del circuito. Lo hice dos veces seguidas en la misma pregunta. Está
contado en la sección~\ref{sec:metodo} porque enseña más que el resultado.

% =====================================================================
\section{El estado del nodo}

\subsection{Qué funciona hoy}

\begin{itemize}
  \item \textbf{Autocalibración desde el peor caso.} Con las cuatro referencias
        en cero y la cadena contra el riel, el nodo la centra sin ayuda externa.
  \item \textbf{Un solo firmware.} El proyecto de autotest dejó de ser una copia
        del de campo: cada fuente compartida existe una vez y la otra la incluye.
        El grabador \emph{se niega a programar} si alguien vuelve a copiar un
        archivo.
  \item \textbf{Batería de ensayos armada} para la sesión de cambio de
        resistencias, con las trampas de la semana metidas como guardas.
\end{itemize}

\subsection{Números vigentes}

\begin{center}
\begin{tabular}{l r p{4.6cm}}
\toprule
\textbf{magnitud} & \textbf{valor} & \textbf{de dónde sale} \\
\midrule
Escala del banco            & 19,9157 & mV reales por mV de banco; dos puntos con tester \\
$V_{ref}$                   & 2,415 V                            & tester \\
$V_{dda}$                   & 4,826 V                            & tester \\
Ventana observable          & 880,4 a 1122,7 & mV de banco; recta anclada en 0 V y $V_{dda}$ \\
$\tau$ de la planta         & 29,5 s                             & cuatro caminos independientes \\
Paso del ADDER sobre el LP  & 61,4 & mV reales por código; barrido en lazo abierto \\
Código que centra la cadena & $-191$                             & barrido en lazo abierto \\
Duración de una calibración & 579 s                              & medido, dos veces \\
Error que consigue          & 34 mV                              & medido, dos veces \\
Máximo PGAout estable a ×50 & ×1                                 & tres corridas \\
\bottomrule
\end{tabular}
\end{center}

% =====================================================================
\section{El hallazgo principal: el firmware se detenía antes del codo}

\begin{figure}[h!]
\centering
\includegraphics[width=\textwidth]{../../lab/planta/figuras/adder_curva_20260906.png}
\caption{Barrido del IDAC del ADDER en lazo abierto, mirando el tap del filtro
pasabajos. La meseta plana de la derecha es el LP saturado: no transmite nada,
haga lo que haga la etapa de arriba. La línea roja es donde el firmware se
detenía. El codo está inmediatamente después.}
\label{fig:adder}
\end{figure}

El filtro pasabajos, saturado, \textbf{no responde en absoluto} hasta que la
contribución del ADDER cruza un umbral cerca del código $-135$. Antes de ese
punto la etapa de arriba puede recorrer volts sin que se vea nada aguas abajo.

\begin{center}
\begin{tabular}{r r l}
\toprule
\textbf{código del ADDER} & \textbf{tap del LP [mV banco]} & \\
\midrule
0 a $-128$ & 748,0 & plano; \textcolor{rojo}{aquí cortaba el firmware} \\
$-144$     & 765,6 & primer movimiento \\
$-160$     & 874,1 & \\
$-176$     & 948,3 & entra en la ventana observable \\
$\mathbf{-191}$ & \textbf{1001,8} & \textbf{$V_{ref}$, cadena centrada} \\
$-240$     & 1122,9 & riel superior \\
\bottomrule
\end{tabular}
\end{center}

El límite del firmware era \ruta{CAL\_DAC\_MAX\_CHANGE\_GEO\_SUM}~$=128$, y era
\emph{simétrico}. El propio comentario del código ya decía que con eso se perdía
el tramo $-178\ldots-128$ ``que sí es útil'', y anotaba los límites asimétricos
como mejora pendiente. Ese tramo perdido era exactamente el que hacía falta.

\subsection{La comprobación que lo separó de la explicación cómoda}

\begin{figure}[h!]
\centering
\includegraphics[width=\textwidth]{../../lab/planta/figuras/rescate_comparado.png}
\caption{A la izquierda, moviendo el ADDER hasta el límite del firmware: su
propio tap recorre 1,73 V y el del LP no se mueve. A la derecha, moviendo el LP
en todo su recorrido: tampoco. Con estos dos solos la conclusión ``el camino
está cortado'' parece inevitable, y es falsa.}
\label{fig:rescate}
\end{figure}

% =====================================================================
\section{La autocalibración}

\begin{figure}[h!]
\centering
\includegraphics[width=\textwidth]{../../lab/planta/figuras/calibracion_antes_despues.png}
\caption{Los cuatro taps antes y después de \texttt{cal}, sin computadora. El
tap del PGA queda donde estaba porque es el offset de entrada amplificado ×50,
que se aceptó explícitamente: no está en el riel, le sobran 1,5 V hasta masa.}
\label{fig:cal}
\end{figure}

La trayectoria muestra el diseño funcionando en dos fases. El ADDER cruza
\emph{a ciegas} los $\sim$190 códigos que separan el reposo del punto útil
---donde ninguna medida significa nada--- y recién cuando el tap entra en la ventana
observable el lazo cerrado toma el control y el LP hace el ajuste fino.

\begin{center}
\begin{tabular}{l r l}
\toprule
\textbf{etapa} & \textbf{medida del lazo} & \\
\midrule
ADDER, al entrar  & 748,5  & fuera de la ventana: no se puede realimentar \\
ADDER, al salir   & 1001,4 & el rescate en lazo abierto lo trajo a $V_{ref}$ +35 mV \\
LP, paso 1        & 1002,2 & +11 mV \\
LP, paso 2        & \textbf{999,9} & \textbf{$-$35 mV} \\
\bottomrule
\end{tabular}
\end{center}

\subsection{Las siete correcciones que hicieron falta}

Ninguna era de sintonía, y las siete eran silenciosas: no daban error y
devolvían números de aspecto normal.

\begin{enumerate}\itemsep2pt
  \item La espera de la planta vivía en una rama de preprocesador que no compila.
  \item Las esperas se contaban en iteraciones del lazo pero se declaran en
        muestras del ADC, que son 78 veces más cortas.
  \item El lazo no esperaba nada después de mover la referencia: medía el pasado.
  \item La curva no lineal del LP se le aplicaba también al ADDER, invirtiéndole
        el signo del lazo.
  \item Recortar la tabla de etapas dejó al instrumento sin poder escribir dos de
        las cuatro referencias, sin avisar.
  \item El límite simétrico frenaba al ADDER antes del codo.
  \item El rescate en lazo abierto quedó antes del control de espera y se
        disparaba cada 30 ms en vez de una vez por $\tau$.
\end{enumerate}

Las tres primeras son la misma confusión repetida: usar el tiempo del ADC
(30 ms) donde hacía falta el de la cadena ($\tau = 29{,}5$ s).

% =====================================================================
\section{¿Vale la autocalibración lo que cuesta?}

\subsection{Sin calibrar, la mayoría de las configuraciones no existen}

\begin{figure}[h!]
\centering
\includegraphics[width=\textwidth]{../../lab/planta/figuras/sin_calibrar_20260906_0856.png}
\caption{Con las cuatro referencias en cero, o sea el nodo sin sistema de
calibración. Las combinaciones en rojo arrancan con al menos una etapa contra el
riel: no capturan nada. No llevan magnitud porque en esa zona la lectura no mide
el tap, y ponerle un número sería inventar precisión.}
\label{fig:sincal}
\end{figure}

\textbf{12 de 14 combinaciones} de ganancia arrancan con una etapa que no
transmite. Las dos que no son ×1\,×1 (peor tap a 128 mV de $V_{ref}$) y ×2\,×1
(426 mV).

La medida se rehízo el 6 de septiembre porque la primera versión no era
reproducible: la misma configuración daba resultados opuestos según de dónde
viniera la cadena. Volver de la saturación tarda mucho más que $2\tau$, así que
con espera fija cada punto heredaba el estado del anterior. Ahora se espera a
que la cadena \emph{deje de moverse}, el orden alterna ganancias altas y bajas
para que un efecto de historia se vea como incoherencia, y la configuración de
campo se repite como control.

\subsection{Y un ajuste de fábrica tampoco alcanzaría}

\begin{figure}[h!]
\centering
\includegraphics[width=\textwidth]{../../lab/planta/figuras/deriva_tendencia.png}
\caption{El punto calibrado, sin volver a tocarlo. La banda verde es lo que la
calibración consigue.}
\label{fig:deriva}
\end{figure}

Sobre \derivaHoras{} h, el archivo contiene \derivaMuestrasTotales{} instantes;
\derivaMuestras{} tienen una lectura válida del tap del filtro pasabajos. Ese
tap recorrió una banda de \textbf{\derivaBanda{} mV}, que son
\textbf{\derivaVeces{} veces} los 34 mV que la calibración consigue.

La auditoría encontró \lecturasTapNulas{} valores nulos entre
\lecturasTapTotales{} lecturas de taps (\lecturasTapNulasPct{}\,\%). El campo
histórico \texttt{lecturas\_fallidas=0} no contaba respuestas vacías; el error
quedó corregido en el registrador para las próximas corridas. Los nulos se
ignoran, no se interpolan.

\paragraph{El registro completo corrige la lectura provisional.}
A las 7,1 h el punto había vuelto cerca del valor inicial y parecía un
vagabundeo acotado. Era demasiado pronto para cerrar esa interpretación: el
registro posterior cruzó el origen y siguió moviéndose. Sobre la ventana
completa, el resumen descriptivo es una tendencia de
\textbf{\derivaTasaSignada{} mV/h}, con un residuo de
\textbf{\derivaResiduo{} mV RMS} alrededor de la recta.

\begin{quote}
arrancó en \derivaArranco{} mV \quad$\rightarrow$\quad bajó hasta
\textbf{\derivaExtremo{} mV} a las \derivaExtremoH{} h \quad$\rightarrow$\quad
y terminó en \derivaAhora{} mV.\\[1mm]
\textbf{Desplazamiento neto en \derivaHoras{} h: \derivaNeto{} mV. Excursión
recorrida: \derivaBanda{} mV.}
\end{quote}

El retorno observado a mitad de la corrida demuestra que el movimiento no es
monótono, pero \emph{no} demuestra que la banda esté cerrada. La descripción
compatible con el registro largo es \textbf{tendencia lenta con vagabundeo
superpuesto}. La recta no se extrapola hasta el riel: se usa sólo para resumir
esta ventana, donde acumula \derivaRectaSignada{} mV.

\paragraph{El argumento de la calibración automática se sostiene.} Un trim de
fábrica no mantiene el tap dentro de los 34 mV que consigue la calibración: la
tendencia consume ese margen en aproximadamente \derivaCome{} h, y el recorrido
total ya es \derivaVeces{} veces mayor. Sólo un lazo que vuelva a medir en el
lugar puede corregir tanto el desplazamiento lento como las excursiones
superpuestas. El mecanismo físico requiere temperatura medida en el propio
banco. Como indicio, la temperatura exterior modelada recorrió
\tempMin--\tempMax{} °C y su correlación contemporánea con el tap fue
$r=\tempCorr{}$ sobre \tempN{} muestras interpoladas. La relación es sugestiva,
pero no reemplaza un sensor junto a la placa ni permite atribuir causalidad.

\subsection{Y calibrar no mete ruido}

Se informa la \textbf{mediana} y no el promedio: el piso típico fue
\ruidoRms{} µV RMS sobre \ruidoN{} medidas válidas
(faltan \ruidoFaltante{} de los \ruidoEsperado{} bloques previstos). Hubo
\ruidoPicos{} por encima del doble del valor típico, hasta \ruidoMax{} µV
(\ruidoVecesMax{} veces la mediana). Esas excitaciones son compatibles con
ruido ambiente ---pasos o tránsito en la lomada---, no con un cambio del piso
de la cadena. Los puntos no se descartan: siguen en los datos y en la figura.
El promedio queda en \ruidoPromedio{} µV, un \ruidoSesgoPromedio{}\,\% por
encima de la mediana, por eso no se usa como estimador del piso.

\begin{figure}[h!]
\centering
\includegraphics[width=\textwidth]{../../lab/planta/figuras/ruido_exp4d.png}
\caption{El ruido del tap del LP con la cadena calibrada y quieta. No crece, y
la componente de red es marginal frente al total.}
\label{fig:ruido}
\end{figure}

% =====================================================================
\section{El ADC usa el 4,8\,\% de su rango}

Dos incógnitas que parecían independientes resultaron ser el mismo número.

El LSB del ADC en su propia entrada son 19,07 µV. Multiplicado por la escala
verificada con tester, \textbf{una cuenta son 379,9 µV de tensión real en el
tap}. Y los bits que eso cuesta:
\[
  \log_2(19{,}9157) = 4{,}32 \quad\Longrightarrow\quad 18 - 4{,}32 = \mathbf{13{,}68\ \text{bits}}
\]
Verificado por el otro camino, contando cuentas: los 4,826 V útiles ocupan
12.703 cuentas, o sea 13,63 bits. Coinciden dentro de 0,05 bits.

\textbf{Los 4,4 bits no se pierden dentro del ADC: se gastan antes de llegar.}
Hay una atenuación de $\approx 19{,}92:1$ entre el tap y la entrada, y el
convertidor termina usando el 4,8\,\% de su rango. Si esa atenuación se redujera
a $2:1$ ---lo justo para meter 5 V en $\pm 2{,}5$ V--- la adquisición ganaría
\textbf{3,3 bits} sin tocar el ADC ni el algoritmo.

Las otras tres configuraciones del ADC, de rango más angosto, \textbf{hoy no se
pueden usar}: los taps se sientan en 0,96--1,12 V a la entrada, que supera su
fondo de escala, y dan la vuelta a valores negativos \emph{con el bit de válido
en~1}. Sólo servirían si antes se quita el pedestal, y ahí serían otros 2,3 bits.

% =====================================================================
\section{Un error de método, contado porque enseña más que el resultado}
\label{sec:metodo}

Concluí dos veces seguidas, con evidencia aparentemente sólida, que el ADDER no
podía rescatar al filtro pasabajos saturado. Barrí su referencia hasta $-128$
códigos y el tap del LP no se movió ni un milivoltio mientras el tap propio del
ADDER recorría 1,73 V (figura~\ref{fig:rescate}). Escribí la conclusión, cambié
el orden de la calibración y lo registré.

$-128$ no era un límite físico: era el límite que el propio firmware le pone a
esa etapa, y yo había elegido el rango del barrido copiando ese número.
Siguiendo hasta $-144$ el tap se movió por primera vez, y en $-191$ la cadena
quedó centrada. \textbf{El umbral estaba dieciséis códigos más allá de donde
corté.}

Lo destapó un dato viejo, no una idea: un archivo guardado decía que el
procedimiento que sí funcionaba usaba el código $-191$, imposible para el
firmware. Sin ese archivo me habría quedado con la conclusión falsa.

\begin{quote}\itshape
Antes de concluir de un barrido que tocó su borde, hay que preguntarse de dónde
salió ese borde. Si salió de una constante del firmware, de un límite o de un
valor copiado de otro lado, el barrido tiene que extenderse hasta el límite del
hardware, y recién ahí se concluye.
\end{quote}

% =====================================================================
\section{Dónde está cada cosa en el repositorio}

\subsection{Documentos}

\begin{longtable}{p{7.6cm} p{8.4cm}}
\toprule
\textbf{archivo} & \textbf{qué tiene} \\
\midrule
\endhead
\ruta{docs/MEDICIONES\_2026-09-05.md} &
  El registro completo, 28 secciones. Las secciones superadas llevan una marca
  que dice cuál las corrige; el cuerpo no se toca, porque el recorrido importa \\
\ruta{docs/MEDICIONES\_2026-09-04.md} &
  Matriz de acople y curvas propias. Ocho secciones marcadas como corregidas \\
\ruta{lab/HANDOFF\_GOAL\_2026-09-05.md} &
  Estado de la tarea. \emph{Empezar por acá} \\
\ruta{lab/PLAN\_RESISTENCIAS\_LUNES.md} &
  El protocolo de la sesión de resistencias, con los comandos que existen \\
\bottomrule
\end{longtable}

\subsection{Firmware}

\begin{longtable}{p{7.6cm} p{8.4cm}}
\toprule
\textbf{ruta} & \textbf{qué es} \\
\midrule
\endhead
\ruta{src/firmware/psoc/AcondicionamientoAnalogico.cydsn/} &
  \textbf{El proyecto de campo. Fuente de verdad de todo} \\
\ruta{\ \ calibration.c} &
  El lazo: rescate en lazo abierto, esperas de planta, PI \\
\ruta{\ \ calibration\_tables*.h} &
  Todos los números por etapa, con su justificación medida al lado \\
\ruta{src/firmware/psoc/AcondicionamientoAnalogicoTest/} &
  El de laboratorio: \emph{incluye} al de campo, no lo copia \\
\ruta{\ \ psoc\_selftest\_primitivas.inc} &
  Lo único propio: las primitivas de medición \\
\ruta{src/firmware/psoc/program\_psoc.ps1} &
  Compila y graba. Verifica la sincronía y se niega a grabar si se rompió \\
\ruta{src/firmware/esp32/.../main\_selftest.cpp} &
  El secuenciador del banco y los comandos del enlace \\
\bottomrule
\end{longtable}

\subsection{Banco y análisis}

\begin{longtable}{p{7.6cm} p{8.4cm}}
\toprule
\textbf{ruta} & \textbf{qué hace} \\
\midrule
\endhead
\ruta{src/interfaces/python/escala\_banco.py} &
  \textbf{La conversión banco$\rightarrow$volts, en un solo lugar}. Y la ventana
  donde la lectura significa algo \\
\ruta{src/interfaces/python/asentamiento.py} &
  Esperar a que la cadena deje de moverse, en vez de un tiempo fijo \\
\ruta{src/interfaces/python/lunes/} &
  \textbf{La batería de ensayos}, seis módulos y un runner \\
\ruta{src/interfaces/python/exp\_deriva.py} &
  EXP4c y EXP4d: deja el nodo calibrado y registra durante horas \\
\ruta{src/interfaces/python/analizar\_deriva.py} &
  Distingue tendencia, vagabundeo y cola de transitorio \\
\ruta{src/interfaces/python/figuras\_semana.py} &
  \textbf{Genera todas las figuras de este informe} desde los JSON crudos \\
\ruta{src/interfaces/python/figura\_sin\_calibrar.py} &
  La figura del argumento de la tesis \\
\bottomrule
\end{longtable}

\subsection{Datos crudos}

Todo en \ruta{lab/planta/}, un JSON por ensayo con fecha en el nombre.

\begin{longtable}{p{7.6cm} p{8.4cm}}
\toprule
\textbf{archivo} & \textbf{qué contiene} \\
\midrule
\endhead
\ruta{rescate\_addr\_largo.json} &
  El barrido completo del ADDER $\rightarrow$ figura~\ref{fig:adder} \\
\ruta{rescate\_lazo\_abierto.json}, \ruta{rescate\_lp.json} &
  Los dos ensayos de rescate $\rightarrow$ figura~\ref{fig:rescate} \\
\ruta{cal\_traza.json} &
  La calibración que funcionó, con su trayectoria $\rightarrow$ figura~\ref{fig:cal} \\
\ruta{sin\_calibrar\_20260906\_0856.json} &
  Las 17 combinaciones de ganancia $\rightarrow$ figura~\ref{fig:sincal} \\
\ruta{deriva\_20260906\_0954.json} &
  EXP4c y EXP4d $\rightarrow$ figuras~\ref{fig:deriva} y~\ref{fig:ruido} \\
\ruta{temperatura\_exterior\_20260906\_20260907.json} &
  Proxy meteorológico horario de Open-Meteo; no es un sensor del banco \\
\ruta{figuras/} &
  Las figuras generadas, en PNG \\
\bottomrule
\end{longtable}

\subsection{Cómo rehacer cualquier cosa de este informe}

\begin{verbatim}
cd src/interfaces/python

python -m lunes.correr_lunes antes      # la huella de la placa
python -m lunes.t1_autoridad 2          # la curva del ADDER
python -m lunes.t3_cal_firmware         # ¿se calibra solo?
python -m lunes.t2_sin_calibrar         # las combinaciones de ganancia
python exp_deriva.py --horas 24         # EXP4c + EXP4d, corte limpio

python descargar_temperatura_deriva.py  # proxy exterior, no sensor del banco
python figuras_semana.py                # regenera todas las figuras
python analizar_deriva.py               # el veredicto sobre la deriva
\end{verbatim}

% =====================================================================
\section{Lo que queda}

\begin{longtable}{p{1.2cm} p{14.8cm}}
\toprule
& \\
\endhead
\textbf{1.} & \textbf{La sesión de resistencias.} Bajar el filtro pasabajos a
1,8 k$\Omega$: gana 8,3× de resolución en el ajuste fino y sigue cubriendo
10,4 pasos del actuador grueso, así que por el lado de los huecos es seguro.
La batería está armada. \\[2mm]
\textbf{2.} & \textbf{El cero de la referencia del ADDER está corrido 358 mV.}
El 75\,\% de su recorrido se gasta en llegar al punto de operación. No es la
resistencia en serie ---su alcance en volts es el correcto--- sino el punto de
reposo. Es una pregunta de esquemático. \\[2mm]
\textbf{3.} & \textbf{Dónde vive la atenuación de 20:1} entre el tap y el ADC.
Vale 3,3 bits de rango dinámico, más otros 2,3 si además se quita el pedestal. \\[2mm]
\textbf{4.} & \textbf{Que el nodo informe éxito también desde el peor caso.}
Hoy informa \texttt{ok=1} arrancando cerca y \texttt{ok=0} arrancando del riel,
con el mismo resultado físico en los dos: es contabilidad de pasos, no
calibración. Diagnosticado, y el ensayo T3 lo confirma solo. \\[2mm]
\textbf{5.} & \textbf{Medir temperatura en el propio banco.} El registro largo
ya separa la lectura prematura de 7,1 h de la tendencia posterior. Falta un
sensor junto a la placa para distinguir deriva térmica de otros mecanismos y
decidir si conviene recalibrar por reloj, por temperatura o por ambos. \\
\bottomrule
\end{longtable}

\end{document}
